\documentclass[12pt]{article}
\usepackage[utf8]{inputenc}
\usepackage[T1]{fontenc}
\usepackage[margin=1in]{geometry}
\usepackage{setspace}
\usepackage{xcolor}
\usepackage[colorlinks=true, urlcolor=blue!60!black, linkcolor=blue!60!black]{hyperref}
\setlength{\parindent}{2em}
\setlength{\parskip}{1.25em}
\renewcommand{\baselinestretch}{1.0}
\title{\textit{Crosby v. National Foreign Trade Council}, 530 U.S. 363 (2000)}
\author{Case Brief}
\date{}
\begin{document}
\maketitle

% Case Brief: Crosby v. National Foreign Trade Council, 530 U.S. 363 (2000)
% Seth J. Chandler's Constitutional Law Briefing Book
% Chapter: Federalism and State Power

\section{Case Metadata}

\begin{description}
    \item[Full Case Name:] \textit{Stephen P.\ Crosby, Secretary of Administration and Finance of Massachusetts, et al.\ v.\ National Foreign Trade Council}
    \item[Citation:] 530 U.S.\ 363 (2000) (No.\ 99--474)
    \item[Court:] Supreme Court of the United States
    \item[Dates:] Argued March 22, 2000; Decided June 19, 2000
    \item[Lower Court Proceedings:] \textit{National Foreign Trade Council v.\ Baker}, 26 F.\ Supp.\ 2d 287 (D.\ Mass.\ 1998) (permanently enjoining enforcement of the Massachusetts Burma Law); affirmed sub nom.\ \textit{National Foreign Trade Council v.\ Natsios}, 181 F.\ 3d 38 (1st Cir.\ 1999); cert.\ granted, 528 U.S.\ 1018 (1999) (caption changed twice as Massachusetts Secretaries of Administration and Finance succeeded one another: Charles D.\ Baker, Andrew S.\ Natsios, and Stephen P.\ Crosby)
    \item[Disposition:] Affirmed (judgment of the First Circuit affirmed solely on Supremacy Clause obstacle preemption grounds)
    \item[Vote:] 9--0 (Unanimous in judgment)
    \item[Majority Opinion:] Justice David H.\ Souter, joined by Chief Justice William H.\ Rehnquist and Justices John Paul Stevens, Sandra Day O'Connor, Anthony M.\ Kennedy, Ruth Bader Ginsburg, and Stephen G.\ Breyer (7 Justices)
    \item[Concurring Opinion:] Justice Antonin Scalia, concurring in the judgment, joined by Justice Clarence Thomas (2 Justices)
    \item[Dissenting Opinions:] None
    \item[Constitutional Provisions \& Statutes:] U.S.\ Const.\ art.\ VI, cl.\ 2 (Supremacy Clause); Foreign Operations, Export Financing, and Related Programs Appropriations Act, 1997, \S~570, Pub.\ L.\ No.\ 104--208, 110 Stat.\ 3009--166 to 3009--167 (1996) (Federal Burma Act); Exec.\ Order No.\ 13047, 62 Fed.\ Reg.\ 28,301 (May 20, 1997); 31 C.F.R.\ pt.\ 537 (1999); 1996 Mass.\ Acts 239, ch.\ 130, codified at Mass.\ Gen.\ Laws ch.\ 7, \S\S~22G--22M (Massachusetts Burma Law)
    \item[Primary Opinion Sources:] \href{https://supreme.justia.com/cases/federal/us/530/363/}{Justia U.S.\ Supreme Court Center (530 U.S.\ 363)}; \href{https://www.law.cornell.edu/supct/html/99-474.ZO.html}{Cornell LII (Majority Opinion)}; \href{https://www.law.cornell.edu/supct/html/99-474.ZC.html}{Cornell LII (Scalia, J., Concurring in Judgment)}; \href{https://www.law.cornell.edu/supct/html/99-474.ZS.html}{Cornell LII (Syllabus)}; \href{https://tile.loc.gov/storage-services/service/ll/usrep/usrep530/usrep530363/usrep530363.pdf}{Library of Congress U.S.\ Reports PDF}; \href{https://www.oyez.org/cases/1999/99-474}{Oyez Project}
\end{description}

\section{Facts and Procedural History}

\subsection{The Massachusetts Burma Law}
In June 1996, reacting to severe human rights abuses by the military junta ruling Burma (which the regime renamed Myanmar in 1989, though the parties, statutes, and Court used ``Burma'' without expressing a political view; \href{https://supreme.justia.com/cases/federal/us/530/363/}{530 U.S.\ at 366 n.\ 1}), the Commonwealth of Massachusetts enacted ``An Act Regulating State Contracts with Companies Doing Business with or in Burma (Myanmar),'' 1996 Mass.\ Acts 239, ch.\ 130, codified at Mass.\ Gen.\ Laws ch.\ 7, \S\S~22G--22M (\href{https://www.law.cornell.edu/supct/html/99-474.ZS.html}{Cornell LII Syllabus}). Modeled on state and municipal selective-purchasing statutes deployed against apartheid-era South Africa in the 1980s, the Massachusetts Burma Law barred state agencies and authorities from buying goods or services from any person or company listed on a state-maintained ``restricted purchase list'' of entities ``doing business with Burma.''

The state statute operated through a rigid procurement penalty and sweeping definitions:
\begin{itemize}
    \item \textbf{Ten Percent Bid Penalty:} Listed companies were effectively disqualified from state contracts unless they submitted a bid at least 10\% lower than the lowest non-listed bidder, or fell within narrow statutory exceptions (such as medical supplies with no substitute, or procurements certified as essential to avoid irreparable harm where applying the restriction would eliminate the only bid or leave inadequate competition) (\href{http://scholarship.law.cornell.edu/cgi/viewcontent.cgi?article=1470&context=cilj}{Cornell Int'l L.J.}).
    \item \textbf{Broad Entity and Activity Coverage:} The law applied equally to domestic and foreign corporations, swept in parent companies and subsidiaries based on ``upstream'' or pre-existing investments and affiliates, and covered sales of goods or services to the Burmese government as well as the promotion of Burmese timber, gems, oil, or gas. The state's Operational Services Division updated the restricted list at least quarterly; companies could challenge inclusion only by submitting an affidavit certifying they did no business with Burma (\href{https://supreme.justia.com/cases/federal/us/530/363/}{530 U.S.\ at 368--369}).
    \item \textbf{Perpetual Duration:} The statute contained no sunset clause, presidential waiver, or termination mechanism tied to diplomatic developments (\href{https://www.law.cornell.edu/supct/html/99-474.ZS.html}{Cornell LII Syllabus}).
\end{itemize}

\subsection{The Federal Burma Act and Executive Implementation}
Three months after Massachusetts acted, in September 1996, Congress enacted \S~570 of the Foreign Operations, Export Financing, and Related Programs Appropriations Act, 1997, Pub.\ L.\ No.\ 104--208, 110 Stat.\ 3009--166 to 3009--167. Congress crafted a calibrated, five-part federal sanctions architecture:
\begin{enumerate}
    \item \textbf{Immediate Mandatory Sanctions:} Imposed targeted immediate restrictions on most bilateral assistance to Burma, required U.S.\ representatives at international financial institutions to oppose loans to Burma, and restricted entry visas for Burmese government officials (\S~570(a)).
    \item \textbf{Conditional Authority to Ban ``New Investment'':} Authorized---and conditioned---the President to prohibit ``new investment'' in Burma by ``United States persons'' if the President determined and certified to Congress that the Burmese regime had physically harmed, rearrested, or exiled democracy leader Aung San Suu Kyi, or committed large-scale repression or violence against the democratic opposition (\S~570(b)). Pre-existing investments, non-U.S.\ persons, and general trade in goods were deliberately excluded.
    \item \textbf{Multilateral Diplomatic Directive:} Directed the President to cooperate with the Association of Southeast Asian Nations (ASEAN) and other countries to develop a ``comprehensive, multilateral strategy'' to improve human rights and restore democracy in Burma, and to report periodically to Congress (\S~570(c)--(d)).
    \item \textbf{Presidential Waiver and Termination Authority:} Empowered the President to suspend or waive any sanction temporarily or permanently upon certifying that doing so was ``in the national security interests of the United States'' (\S~570(e)), and to terminate sanctions upon certifying measurable progress in human rights and democracy.
\end{enumerate}

On May 20, 1997, President Bill Clinton issued Executive Order No.\ 13047, certifying large-scale repression by the Burmese government, declaring a national emergency, and prohibiting new investment in Burma by U.S.\ persons (and facilitation or evasion thereof), while incorporating the statutory exemptions (\href{https://supreme.justia.com/cases/federal/us/530/363/}{530 U.S.\ at 369--370}). On May 21, 1998, the Treasury Department promulgated implementing Burmese Sanctions Regulations, 31 C.F.R.\ pt.\ 537 (\href{https://supreme.justia.com/cases/federal/us/530/363/}{530 U.S.\ at 370 n.\ 3}).

Meanwhile, the Massachusetts Burma Law provoked sharp international backlash. The European Union and Japan lodged formal diplomatic protests with the U.S.\ government and initiated dispute-settlement proceedings against the United States before the World Trade Organization (WTO), alleging violations of the WTO Agreement on Government Procurement (\href{https://www.law.cornell.edu/supct/html/99-474.ZS.html}{Cornell LII Syllabus}).

\subsection{Litigation History}
In April 1998, the National Foreign Trade Council (NFTC), a nonprofit trade association representing American companies engaged in foreign commerce---34 of whose members appeared on Massachusetts's restricted purchase list (\href{https://www.oyez.org/cases/1999/99-474}{Oyez})---sued Massachusetts officials in federal district court. The NFTC sought declaratory and injunctive relief on three constitutional grounds: (1) intrusion upon the federal government's exclusive foreign affairs power under \textit{Zschernig v.\ Miller}, 389 U.S.\ 429 (1968); (2) violation of the dormant Foreign Commerce Clause, U.S.\ Const.\ art.\ I, \S~8, cl.\ 3; and (3) Supremacy Clause preemption by the Federal Burma Act.

The U.S.\ District Court for the District of Massachusetts permanently enjoined the statute, holding that it unconstitutionally impinged on the federal government's exclusive authority to regulate foreign affairs (\textit{Nat'l Foreign Trade Council v.\ Baker}, 26 F.\ Supp.\ 2d 287, 291 (D.\ Mass.\ 1998)). The First Circuit affirmed on all three independent grounds: dormant foreign affairs preemption under \textit{Zschernig}, the dormant Foreign Commerce Clause, and statutory conflict preemption (\textit{Nat'l Foreign Trade Council v.\ Natsios}, 181 F.\ 3d 38, 45 (1st Cir.\ 1999)). The Supreme Court granted certiorari (\href{https://caselaw.findlaw.com/court/us-supreme-court/530/363.html}{528 U.S.\ 1018 (1999)}).

\section{Issue(s) Presented}

Whether the Massachusetts Burma Law, which restricts state agencies from purchasing goods or services from companies doing business with Burma, is impliedly preempted---and therefore unconstitutional under the Supremacy Clause---because it stands as an obstacle to the accomplishment and execution of the full purposes and objectives of Congress in the Federal Burma Act.

\textit{Note on Issues Reserved:} Because the Court resolved the case solely on statutory obstacle preemption, it expressly declined to decide (1) whether a presumption against preemption applies in the foreign affairs context, (2) whether Congress occupied the field (field preemption), (3) whether the law violated the dormant foreign affairs power under \textit{Zschernig v.\ Miller}, or (4) whether the law violated the dormant Foreign Commerce Clause (or whether Massachusetts was shielded by a market-participant exception) (\href{https://www.law.cornell.edu/supct/html/99-474.ZO.html}{530 U.S.\ at 374 n.\ 8}).

\section{Holding}

\textbf{Yes.} Writing for a seven-Justice majority (with two Justices concurring in the judgment), Justice Souter held that the Massachusetts Burma Law is impliedly preempted under the Supremacy Clause because it stands as an obstacle to the accomplishment and execution of Congress's full statutory objectives in the Federal Burma Act (\href{https://www.law.cornell.edu/supct/html/99-474.ZO.html}{530 U.S.\ at 366, 388}). Specifically, the state law undermined the intended purpose and natural effect of three core features of the federal statute: (1) its delegation of broad, flexible discretion to the President to control and lift economic sanctions; (2) its deliberate calibration limiting sanctions to U.S.\ persons and ``new investment''; and (3) its directive to the President to speak for the United States in crafting a comprehensive, multilateral diplomatic strategy.

\section{Reasoning}

\subsection{Preemption Framework and Rejection of the ``Spending Power'' Defense}
Justice Souter began from settled Supremacy Clause principles: even absent an express preemption clause, state law must yield when Congress intends to occupy a field, when dual compliance is physically impossible, or when ``under the circumstances of [a] particular case, [the challenged state law] stands as an obstacle to the accomplishment and execution of the full purposes and objectives of Congress'' (\textit{Hines v.\ Davidowitz}, 312 U.S.\ 52, 67 (1941); \href{https://caselaw.findlaw.com/court/us-supreme-court/530/363.html}{530 U.S.\ at 372--373}). Judging whether a state law poses a sufficient obstacle requires examining the federal statute as a whole to identify its purpose and ``natural effect.''

Crucially, the Court dispatched Massachusetts's threshold defense that it was merely acting as a ``market participant'' exercising its proprietary spending power rather than its regulatory police power. Relying on \textit{Wisconsin Department of Industry v.\ Gould Inc.}, 475 U.S.\ 282, 287--289 (1986), the Court reiterated that a state statutory scheme does not escape Supremacy Clause preemption merely because the state chooses to wield its purchasing power rather than direct regulation to enforce a policy that conflicts with a federal statute (\href{https://tile.loc.gov/storage-services/service/ll/usrep/usrep530/usrep530363/usrep530363.pdf}{530 U.S.\ at 373 n.\ 7}). The Court also assumed \textit{arguendo} that a presumption against preemption might apply, finding that the Massachusetts statute failed even under that forgiving assumption (\href{https://www.law.cornell.edu/supct/html/99-474.ZO.html}{530 U.S.\ at 374 n.\ 8}).

\subsection{Pillar One: Obstruction of Presidential Discretion and ``Bargaining Chips''}
First, the Court emphasized that the Federal Burma Act delegated sweeping, flexible discretion to the President---to impose conditional sanctions, to suspend them for national security reasons (\S~570(e)), and to terminate them upon democratic progress. Invoking Justice Jackson's first tier in \textit{Youngstown Sheet \& Tube Co.\ v.\ Sawyer}, 343 U.S.\ 579, 635 (1952), Justice Souter stressed that when the President acts pursuant to an express congressional authorization in foreign affairs, his authority is at its maximum (\href{https://www.law.cornell.edu/supct/html/99-474.ZS.html}{530 U.S.\ at 375--376}).

This ``plenitude of Executive authority'' controlled the preemption inquiry. Congress empowered the President not merely to make a moral statement, but to achieve a concrete diplomatic result. Massachusetts's law, by contrast, operated immediately, rigidly, and perpetually, with no mechanism to yield to presidential waivers. Consequently, Massachusetts stripped the President of full control over the coercive power of the American economy. Borrowing the metaphor from \textit{Dames \& Moore v.\ Regan}, 453 U.S.\ 654, 673 (1981), the Court explained that the state statute ``reduces the value of the [bargaining] chips created by the federal statute,'' leaving the President with less economic and diplomatic leverage to offer Burma in exchange for reform (\href{https://caselaw.findlaw.com/court/us-supreme-court/530/363.html}{530 U.S.\ at 376--377}).

\subsection{Pillar Two: Undermining Congress's Calibration of Economic Pressure}
Second, the Massachusetts statute directly contradicted Congress's deliberate decision to restrict the scope of economic pressure against Burma. The state law penalized (a) companies with pre-existing investments or affiliates, (b) ordinary trade in goods and services rather than solely ``new investment,'' and (c) foreign corporations in addition to U.S.\ persons---all of which Congress explicitly excluded or exempted in \S~570 (\href{https://supreme.justia.com/cases/federal/us/530/363/}{530 U.S.\ at 378--379}).

Massachusetts argued that no conflict existed because both statutes pursued the same ultimate goal (promoting human rights in Burma) and private companies could theoretically comply with both by ceasing all Burmese dealings. The Court firmly rejected this argument: ``[t]he fact of a common end hardly neutralizes conflicting means'' (\textit{Gade v.\ National Solid Wastes Management Ass'n}, 505 U.S.\ 88, 103 (1992)). Sanctions regimes are calibrated instruments---designed not only to forbid what they prohibit, but also to permit what they leave unregulated. By adding harsher, broader state penalties on top of federal sanctions, Massachusetts undermined ``the congressional calibration of force'' (\href{https://supreme.justia.com/cases/federal/us/530/363/}{530 U.S.\ at 379--380}).

\subsection{Pillar Three: Interference with Multilateral Diplomacy and ``One Voice''}
Third, \S~570(c) directed the President to work with ASEAN and other nations to craft a comprehensive, multilateral Burma strategy. The Court concluded that the Massachusetts statute actively frustrated that diplomatic mandate. Foreign allies---including the European Union and Japan---had protested the state law to the federal government and initiated formal WTO dispute proceedings, forcing U.S.\ diplomats to spend diplomatic capital defending Massachusetts rather than building a united coalition against Burma (\href{https://www.law.cornell.edu/supct/html/99-474.ZS.html}{530 U.S.\ at 381--384}).

The Court treated the formal protests of foreign sovereigns and the consistent representations of the Executive Branch as ``competent and direct evidence'' that the state law impaired the President's ability to ``speak for the United States among the world's nations'' (\href{https://www.law.cornell.edu/supct/html/99-474.ZS.html}{530 U.S.\ at 381, 385--386}). The Court distinguished \textit{Barclays Bank PLC v.\ Franchise Tax Board of California}, 512 U.S.\ 298 (1994), where foreign protests had not invalidated California's unitary corporate tax: in \textit{Barclays}, Congress had taken no affirmative action and had repeatedly declined to preempt state tax methods, leaving the case to the dormant Foreign Commerce Clause; in \textit{Crosby}, by contrast, Congress had enacted an explicit statutory directive tasking the President with multilateral diplomacy.

\subsection{Rejection of Congressional Silence as Implicit Permission}
Finally, Massachusetts contended that because Congress knew about the Massachusetts Burma Law when it enacted \S~570 in September 1996---and knew of similar state anti-apartheid laws during the 1980s---its failure to include an express preemption clause signaled tacit approval. The Court rejected this inference as ``unwarranted'' because ``the silence of Congress is ambiguous'' (\href{https://www.law.cornell.edu/supct/html/99-474.ZO.html}{530 U.S.\ at 387--388}). Congress may omit an express preemption clause precisely because it relies on courts to apply settled implied preemption doctrine, and the Court had never adjudicated the validity of 1980s state South Africa divestment laws.

\section{Concurring and Dissenting Opinions}

\subsection{Justice Scalia, Concurring in the Judgment (Joined by Justice Thomas)}
Justice Scalia agreed completely with the Court's judgment and its substantive preemption conclusion, but refused to join Justice Souter's opinion because of its extensive reliance on legislative history---statements by individual bill sponsors, floor remarks, Executive Branch letters to committees, and unpassed competing bills (such as H.R.\ 3540)---spread across roughly a dozen footnotes (\href{https://www.law.cornell.edu/supct/html/99-474.ZC.html}{530 U.S.\ at 388--391}).

Scalia's concurrence is a classic textualist manifesto organized around a rhythmic refrain: on every single point of preemption (presidential flexibility, calibrated scope, multilateral strategy, and diplomatic friction), the statutory text and record made the conclusion ``perfectly obvious on the face of this statute'' (\href{https://www.law.cornell.edu/supct/html/99-474.ZC.html}{530 U.S.\ at 388--390}). Invoking legislative history was therefore both methodologically illegitimate (since floor speeches delivered to an empty chamber and unpassed bills do not reflect the intent of a majority of both Houses) and completely unnecessary. In a celebrated literary flourish, Scalia compared the majority's gratuitous citation of legislative history to St.\ Augustine's famous remorse in \textit{The Confessions} at stealing pears when he was not even hungry, ``just for the devil of it'' (\href{https://www.law.cornell.edu/supct/html/99-474.ZC.html}{530 U.S.\ at 390--391}). He also warned of practical harm: devoting roughly one-tenth of the opinion to legislative-history footnotes signals to appellate counsel that they must exhaustively research and brief legislative history even when statutory text is clear, needlessly inflating litigation costs.

\subsection{Dissenting Opinions}
None. All nine Justices agreed that the Massachusetts Burma Law was preempted by the Federal Burma Act.

\section{Key Quotations}

\begin{enumerate}
    \item \textbf{Three-Part Obstacle Preemption Holding (Justice Souter, Majority Opinion):}
    \begin{quote}
        ``We find that the state law undermines the intended purpose and `natural effect' of at least three provisions of the federal Act, that is, its delegation of effective discretion to the President to control economic sanctions against Burma, its limitation of sanctions solely to United States persons and new investment, and its directive to the President to proceed diplomatically in developing a comprehensive, multilateral strategy towards Burma.''
    \end{quote}
    --- 530 U.S.\ at 373--374 (\href{https://caselaw.findlaw.com/court/us-supreme-court/530/363.html}{FindLaw} / \href{https://www.law.cornell.edu/supct/html/99-474.ZO.html}{Cornell LII}).

    \item \textbf{Plenitude of Executive Authority and Diplomatic ``Bargaining Chips'' (Justice Souter, Majority Opinion):}
    \begin{quote}
        ``Within the sphere defined by Congress, then, the statute has placed the President in a position with as much discretion to exercise economic leverage against Burma, with an eye toward national security, as our law will admit. And it is just this plenitude of Executive authority that we think controls the issue of preemption here. The President has been given this authority not merely to make a political statement but to achieve a political result, and the fullness of his authority shows the importance in the congressional mind of reaching that result. It is simply implausible that Congress would have gone to such lengths to empower the President if it had been willing to compromise his effectiveness by deference to every provision of state statute or local ordinance that might, if enforced, blunt the consequences of discretionary Presidential action. \dots\ Quite simply, if the Massachusetts law is enforceable the President has less to offer and less economic and diplomatic leverage as a consequence. In \textit{Dames \& Moore v.\ Regan}, 453 U.S.\ 654 (1981), we used the metaphor of the bargaining chip to describe the President's control of funds valuable to a hostile country, \textit{id.}, at 673; here, the state Act reduces the value of the chips created by the federal statute.''
    \end{quote}
    --- 530 U.S.\ at 375--377 (\href{https://caselaw.findlaw.com/court/us-supreme-court/530/363.html}{FindLaw} / \href{https://supreme.justia.com/cases/federal/us/530/363/}{Justia}).

    \item \textbf{Common Ends, Conflicting Means, and Calibration of Force (Justice Souter, Majority Opinion):}
    \begin{quote}
        ``The fact of a common end hardly neutralizes conflicting means, see \textit{Gade v.\ National Solid Wastes Management Assn.}, 505 U.S.\ 88, 103 (1992), and the fact that some companies may be able to comply with both sets of sanctions does not mean that the state Act is not at odds with achievement of the federal decision about the right degree of pressure to employ. \dots\ Sanctions are drawn not only to bar what they prohibit but to allow what they permit, and the inconsistency of sanctions here undermines the congressional calibration of force.''
    \end{quote}
    --- 530 U.S.\ at 379--380 (\href{https://supreme.justia.com/cases/federal/us/530/363/}{Justia}).

    \item \textbf{Ambiguity of Congressional Silence (Justice Souter, Majority Opinion):}
    \begin{quote}
        ``A failure to provide for preemption expressly may reflect nothing more than the settled character of implied preemption doctrine that courts will dependably apply, and in any event, the existence of conflict cognizable under the Supremacy Clause does not depend on express congressional recognition that federal and state law may conflict, \textit{Hines}, 312 U.S., at 67. The State's inference of congressional intent is unwarranted here, therefore, simply because the silence of Congress is ambiguous.''
    \end{quote}
    --- 530 U.S.\ at 387--388 (\href{https://www.law.cornell.edu/supct/html/99-474.ZO.html}{Cornell LII}).

    \item \textbf{Textualist Critique of Gratuitous Legislative History (Justice Scalia, Concurring in the Judgment):}
    \begin{quote}
        ``The only reliable indication of that intent---the only thing we know for sure can be attributed to all of them---is the words of the bill that they voted to make law. In a way, using unreliable legislative history to confirm what the statute plainly says anyway (or what the record plainly shows) is less objectionable since, after all, it has absolutely no effect upon the outcome. But in a way, this utter lack of necessity makes it even worse---calling to mind St.\ Augustine's enormous remorse at stealing pears when he was not even hungry, and just for the devil of it (`not seeking aught through the shame, but the shame itself!').''
    \end{quote}
    --- 530 U.S.\ at 390--391 (\href{https://www.law.cornell.edu/supct/html/99-474.ZC.html}{Cornell LII}).
\end{enumerate}

\section{Doctrinal Significance and Modern Legacy (Through 2026)}

\textit{Crosby} is the modern anchor of statutory obstacle preemption in foreign affairs. By ruling 9--0 on statutory grounds rather than reaching \textit{Zschernig v.\ Miller} or the dormant Foreign Commerce Clause, the Court channeled foreign-affairs federalism through Congress, preserving legislative authority to permit or bar state sanctions.

\begin{itemize}
    \item \textbf{Subsequent Supreme Court Milestones:} Three years later, \href{https://supreme.justia.com/cases/federal/us/539/396/}{\textit{American Insurance Ass'n v.\ Garamendi}, 539 U.S.\ 396 (2003)}, extended \textit{Crosby}'s conflict-preemption reasoning from federal statutes to sole executive agreements, striking down California's Holocaust-insurance disclosure law 5--4. In \href{https://www.everycrsreport.com/reports/R42719.html}{\textit{Arizona v.\ United States}, 567 U.S.\ 387 (2012)}, the Court invoked \textit{Crosby}'s ``calibration of force'' principle to preempt state immigration sanctions obstructing federal enforcement discretion. And \href{https://supreme.justia.com/cases/federal/us/576/1/}{\textit{Zivotofsky v.\ Kerry}, 576 U.S.\ 1 (2015)}, quoted \textit{Crosby} for the rule that the Nation must ``speak \dots\ with one voice'' in foreign recognition.
    \item \textbf{Congressional Overrides and Narrower State Laws:} After \textit{Crosby} led a federal court to invalidate Illinois's Sudan sanctions in \textit{NFTC v.\ Giannoulias}, 523 F.\ Supp.\ 2d 731 (N.D.\ Ill.\ 2007), Congress enacted the Sudan Accountability and Divestment Act of 2007 and CISADA (2010) to authorize tailored state divestment expressly (\href{https://www.everycrsreport.com/reports/RL33948.html}{CRS Report RL33948}). Courts also distinguished \textit{Crosby} where states tracked federal designations without broader penalties, as in \textit{Faculty Senate of FIU v.\ Winn}, 616 F.\ 3d 1206 (11th Cir.\ 2010) (\href{https://www.everycrsreport.com/reports/RL33948.html}{CRS Report RL33948}).
    \item \textbf{Textualist Pushback and 2025--2026 Frontiers:} Recently, textualist Justices have curbed implied obstacle preemption outside core foreign-affairs statutes---requiring conflict grounded in statutory text rather than a ``brooding federal interest'' in \href{https://www.law.cornell.edu/supremecourt/text/16-1275}{\textit{Virginia Uranium, Inc.\ v.\ Warren}, 587 U.S.\ 761 (2019)}, and questioning ``purposes and objectives'' preemption in \href{https://www.scotusblog.com/2020/03/opinion-analysis-divided-court-permits-state-identity-theft-prosecution-of-noncitizens-in-the-employment-process/}{\textit{Kansas v.\ Garcia}, 589 U.S.\ 191 (2020)}. In \href{https://media.ca11.uscourts.gov/opinions/pub/files/202312737.pdf}{\textit{Shen v.\ Simpson} (11th Cir.\ Nov.\ 4, 2025)}, the Eleventh Circuit distinguished \textit{Crosby} in denying a preliminary injunction against Florida's foreign-adversary land-ownership law (S.B.\ 264) under CFIUS. Through 2026, defendants in state climate-tort suits have likewise invoked \textit{Crosby} to argue foreign-affairs preemption, though courts and scholars largely resist extending its statute-specific holding to domestic tort remedies (\href{https://lawreview.uchicago.edu/online-archive/polluter-pays-principle-and-its-enemies-corporate-fraud-weaponizing-preemption-and}{U.\ Chi.\ L.\ Rev.\ Online (2026)}).
\end{itemize}

\section{Pedagogical Critiques and Exam Hypotheticals}

\subsection{Pedagogical Critiques}
\begin{itemize}
    \item \textbf{Market-Participant Asymmetry:} While states enjoy a ``market participant'' shield under the domestic dormant Commerce Clause, \textit{Crosby} (via \textit{Gould}) confirms that spending-power exercises enjoy no immunity from Supremacy Clause preemption.
    \item \textbf{Elasticity and Foreign Protests:} Treating any partial federal sanction as a ``calibration of force'' invites expansive obstacle preemption, while crediting foreign diplomatic protests risks giving foreign sovereigns a de facto veto over subnational policy.
\end{itemize}

\subsection{Exam Hypotheticals}
\begin{itemize}
    \item \textbf{Classic \textit{Crosby} Preemption:} Congress authorizes the President to ban new U.S.\ energy investment in Freedonia, subject to a national-security waiver. New York bars state agencies from contracting with \textit{any} domestic or foreign firm doing any business in Freedonia, without a waiver, sparking EU protests. \textit{Result:} Preempted under \textit{Crosby}---the state spending defense fails, and the law blunts presidential waiver leverage, overrides Congress's calibrated scope, and impedes multilateral diplomacy.
    \item \textbf{``Opposite Side'' Contrast (Surviving \textit{Crosby}):} Congress enacts a Freedonia Divestment Act expressly authorizing states to divest public pensions from firms on a Treasury defense list, provided state measures automatically sunset upon federal delisting or presidential waiver. Texas enacts a matching pension-divestment statute. \textit{Result:} Upheld---express congressional authorization overcomes \textit{Crosby}, \textit{Zschernig}, and the dormant Foreign Commerce Clause, while automatic sunsetting and federal-list tracking preserve both presidential ``bargaining chips'' and congressional calibration.
\end{itemize}

\end{document}
